\documentclass[10pt,a4paper]{amsart}
\usepackage[margin=19mm]{geometry}
\usepackage{amsmath,amssymb,mathtools}
\usepackage[scr=rsfs,cal=euler]{mathalfa}
\usepackage{xfrac}
\usepackage[unicode,hidelinks,bookmarks=false]{hyperref}
\newcommand{\ie}{i.e.\ }
\newcommand{\cf}{cf.\ }
\newcommand{\resp}{respectively\ }
\newcommand{\Subsection}{\subsection}
\providecommand{\nobreakdash}{\nobreak\hbox{-}}
\setlength{\emergencystretch}{2em}
\multlinegap0pt
\usepackage{bm}
\usepackage{bbm}
\usepackage{dsfont}
\usepackage{xspace}
\usepackage{cancel}
\usepackage{mathtools}
\usepackage{amsthm}
\makeatletter
\@ifundefined{lemma}{\newtheorem{lemma}{Lemma}}{}
\@ifundefined{corollary}{\newtheorem{corollary}[lemma]{Corollary}}{}
\@ifundefined{proposition}{\newtheorem{proposition}[lemma]{Proposition}}{}
\makeatother
\newcommand{\id}[1]{\langle #1 \rangle}
\newcommand{\lm}{\lambda}
\title[Cyclic and Constacyclic Codes Over Z4+iZ4]{Cyclic and Constacyclic Codes Over Z4+iZ4}
\author{Miguel Mart\'in, Ekin \"Ozman (Faculty of Science, Engineering - Bernoulli Institute, University of Groningen)}
\date{Source: arXiv:2607.18471v1 (CC BY 4.0)}
\begin{document}
\maketitle

\noindent\emph{Source and licence.} Cyclic and Constacyclic Codes Over Z4+iZ4. Version arXiv:2607.18471v1 (CC BY 4.0), distributed under the Creative Commons Attribution 4.0 International licence, as recorded in the arXiv licence metadata for this exact version and shown on the abstract page https://arxiv.org/abs/2607.18471v1. Full rights notice: licenses/arxiv-2607.18471-CC-BY-4.0.txt.\par\smallskip

\noindent\textit{Editorial scope.} Counted unit: the Proposition whose source label is \texttt{equivalencesR} in the article (source0.tex lines 677--684, source label \texttt{equivalencesR}), delivered together with the article's own complete original proof of it (source0.tex lines 686--706, one \texttt{proof} environment of 21 lines). No mathematics is written, repaired or re-derived: statement and proof are the article's own text line for line. Required context retained uncounted: caseseq (lines 662--670). Every definition, display and label of those lines is retained unchanged. Omitted from this package and not redistributed: 1--661, 671--676, 685--685, 707--1360. Nothing inside a retained excerpt is omitted or altered: every retained body is delivered exactly as the article's lines are written, so no omission receipt of any kind accompanies this package. Citations: the retained excerpts carry no citation command at all, so no bibliography entry of the article is transcribed and no reference is redistributed. Reference adaptation: none was needed and none was made; every reference inside the delivered unit resolves inside it, so no unresolved reference or citation is produced. Printed slips kept literal: no printed word, symbol, hypothesis or number of the counted statement or of its proof was altered. Layout and class: the article's own class is \texttt{amsart} with packages including amssymb, amsthm, amsmath, amsfonts, amscd, amsbsy, upref, makecell, yhmath, adjustbox, hyperref, amsrefs, xy, graphicx; the wrapper is the lane's pinned \texttt{amsart} editorial wrapper at 10pt on A4, which restates the theorem-like environments the retained text needs and adds the article's own macro definitions that the retained text needs (\texttt{\textbackslash id}, \texttt{\textbackslash lm}), with extra wrapper packages (bm, bbm, dsfont, xspace, cancel, mathtools). The delivered numbering is the unit's own. Assets: no retained span contains a figure environment, a graphics-inclusion command, a PDF-page inclusion, a table or a TikZ picture; no image file is copied into this package, the package declares no asset dependency and no image file is redistributed with it. Licence: version arXiv:2607.18471v1 of this article is distributed under the Creative Commons Attribution 4.0 International licence, as recorded in the arXiv licence metadata for this exact version and shown on the abstract page https://arxiv.org/abs/2607.18471v1. A full rights notice is delivered at licenses/arxiv-2607.18471-CC-BY-4.0.txt. Characters: every retained excerpt is pure ASCII, so the delivered engine, XeLaTeX with its default Latin Modern text font, reports no missing character.
\begin{corollary}\label{caseseq}
    Let $n$ be an odd positive integer. Then we have $\lm^n = \lm$ if and only if 
    \begin{itemize}
        \item $\lm \in U_1$ or
        \item $\lm \in U_2$ and $n \equiv 1 \mod 4$.
    \end{itemize}
         Similarly, we have  $\lm^n = -\lm$ if and only if $\lm \in U_2$ and $n \equiv 3 \mod 4$.
   
\end{corollary}

\begin{proposition}\label{equivalencesR}
    Let $n$ be an odd positive integer. For any unit $\lm$ let $\mu = \pm \lm$ such that $\lm = \mu^n$ according to Corollary \ref{caseseq}. Then the following map is a ring isomorphism that is also a monomial transformation:
    \[\eta: R[x]/\langle x^n-\lm \rangle \to R[x]/\langle x^n- 1 \rangle\]
        \[\overline g(x) \mapsto \overline g(\mu x).\]
        

    
\end{proposition}

\begin{proof}
    First we check that $\eta$ is well-defined,

\[\eta(g(x) + h(x)(x^n-\lm) + \langle x^n - \lm \rangle) = g(\mu x) + h(\mu x)(\mu ^n x^{n}-\lm) + \langle x^n -1\rangle =\]
\[g(\mu x) + h(\mu x)(\lm x^n -\lm) + \langle x^n - 1 \rangle =
g(\mu x) + \langle x^n - 1 \rangle = \eta(g(x) + \langle x^n - \lm \rangle)\]

So the choice of representative does not matter.
This is trivially a homomorphism, to argue this function is an isomorphism one easily finds
\[\eta(g(x) + \id{x^n-\lm}) = \id{x^n-1} \leftrightarrow g(\mu x) \in \id{x^n-1} \leftrightarrow g(x) \in \id{x^n-\lm}\]
so the kernel is trivial, and surjectivity follows from the fact that their cardinality is equal.
The map $\eta$ is clearly a monomial transformation, as each term of the polynomial is mapped to a scalar multiple of itself.

    % To see that its a monomial transformation is not hard, each term $\alpha_k x^k$ of $g$ will just be sent to $\lm^{-k}\alpha_k x^k$, so this is a scaling by units of each coefficient.
    % We need to check well-definedness, but first notice its clear that this respects addition and multiplication. Moreover

    % \[x^n-1 + \langle x^n-1 \rangle \mapsto \lambda^{-n}x^n-1 + \langle x^n - \lm \rangle = \lm (x^n-\lambda) + \langle x^n-\lm \rangle = 0 + \langle x^n - \lm \rangle\]

    % With this we essentially get well-definedness as it implies adding $x^n-1$ to the representative we choose does not alter the result. So it is independent of representative.
    % It is a bijection as the map is invertible by $h + \langle x ^n - \lm \rangle \mapsto h(\lm x) + \langle x^n - 1 \rangle$.
\end{proof}

\end{document}
